(1) 接触点滑速 \(u=v-R\omega\)，\(u_0=v_0-R\omega_0\)。当 \(u_0\ne0\)，摩擦 \(F=-\mu_kmg\,\operatorname{sgn}(u_0)\)，其力矩使自转朝消除滑动方向改变；滑动结束前摩擦符号恒定。
(2) 记 \(s=\operatorname{sgn}u_0\)，\(\dot v=-\mu_kg s\)，\(\dot\omega=5\mu_kg s/(2R)\)，所以 \(u(t)=u_0-7\mu_kgst/2\)。首次 \(u=0\) 时 \(t_*=2|u_0|/(7\mu_kg)\)，之后无外水平力可匀速滚动。
(3) 积分得 \(v_*=v_0-2u_0/7\)，\(\omega_*=\omega_0+5u_0/(7R)\)；代入 \(v_*-R\omega_*=0\)。即使初始球体逆向自旋，末态仍满足滚动约束，但末态平动方向取决于 \(v_0,\omega_0\)。
(4) \(\Delta E=m(v_0^2-v_*^2)/2+I(\omega_0^2-\omega_*^2)/2=mu_0^2/7\)。摩擦功率 \(F u=-\mu_kmg|u|\)，积分得损失 \(\mu_kmg|u_0|t_*/2=mu_0^2/7\)。若 \(u_0=0\)，则 \(t_*=0\) 且不耗散。
